\documentclass[11pt,a4paper]{article}
\usepackage[T1]{fontenc}
\usepackage[spanish]{babel}
\usepackage[margin=2.5cm]{geometry}
\usepackage{graphicx}
\usepackage{tabularx}
\usepackage{booktabs}
\usepackage{fvextra}
\RecustomVerbatimEnvironment{verbatim}{Verbatim}{
  frame=single,
  framesep=3mm,
  fontsize=\small,
  breaklines=true,
  breakanywhere=true
}
\usepackage[backend=biber,style=ieee]{biblatex}
\addbibresource{referencias.bib}
\usepackage[hidelinks]{hyperref}

\title{Comunicación entre procesos con sockets TCP y gRPC\\GA-SUM-01 / PE-U1}
\author{Coronel Gomez Khriz Alexander\\Giraldo Cagua Diego Jose\\
Loor Mendoza Byron Daniel\\Silva Triviño John Jairo}
\date{Noviembre de 2026}

\begin{document}
\maketitle
\noindent Repositorio: \url{https://github.com/Danxxx777/GA-SUM-02-PE-U1}

\section{Objetivo e hipótesis}
Comparar la latencia de ida y vuelta y el rendimiento de los dos mecanismos
en tres procesos. Verificar en los eventos registrados que cada recepción
tiene un reloj lógico mayor que su envío correspondiente.

\paragraph{Hipótesis propuesta, pendiente de validación por el equipo.}
En localhost, se espera que TCP con JSON tenga menor latencia mediana para
cargas de 100 bytes en modo secuencial. La implementación usa una conexión
persistente y un delimitador de línea, mientras que el RPC agrega el
procesamiento de HTTP/2 y su contrato binario. La diferencia con 10240 bytes
se evaluará por separado.

\section{Fundamento y estudios relacionados}
TCP ofrece un flujo de bytes entre extremos \cite{rfc9293}. La aplicación
delimita los mensajes JSON para distinguir solicitudes y acuses.
HTTP/2 define tramas y permite multiplexar intercambios sobre una conexión
\cite{rfc9113}. En esta práctica gRPC usa ese protocolo y un contrato
Protocol Buffers.

Niswar et al. midieron tiempo de respuesta y utilización de CPU al comparar
gRPC, REST y GraphQL en microservicios \cite{niswar2024}.
Blinowski y Pełka evaluaron cuatro familias de RPC en ambientes de nube y
analizaron la influencia del tamaño de los argumentos \cite{blinowski2026}.
Di Meglio y Starace estudiaron cómo el framework y el entorno afectan el
rendimiento y el consumo de recursos \cite{dimeglio2024}.
Estos estudios motivan mantener constantes la máquina y la carga para esta
comparación en localhost.

\section{Arquitectura}
Cada proceso puede enviar y recibir. Los nodos A, B y C escuchan en
\texttt{127.0.0.1} en los puertos 5000, 5001 y 5002, respectivamente.
La secuencia principal recorre A $\to$ B, B $\to$ C y C $\to$ A.
\IfFileExists{figuras/arquitectura.pdf}{%
  \includegraphics[width=\linewidth]{figuras/arquitectura.pdf}%
}{\par\textit{Pendiente de incorporar el diagrama actualizado de FORM-TL-01.}}

\section{Implementación de la comunicación mediante sockets TCP}
\label{sec:sockets}

Cada nodo ejecuta un servidor TCP en un hilo de aceptación y un cliente que
recibe órdenes por la entrada estándar. Esto permite que un mismo proceso
envíe solicitudes y atienda las de otros nodos. Los servidores escuchan
exclusivamente en la dirección local \texttt{127.0.0.1}. Los puertos de la
práctica son 5000, 5001 y 5002 para los nodos A, B y C, respectivamente.

El código se divide en tres módulos. \texttt{servidor.py} implementa el
encuadre de los mensajes, su validación y la atención de conexiones;
\texttt{cliente.py} gestiona los envíos, los acuses y la latencia;
\texttt{nodo.py} integra ambos con el reloj y el registro compartidos.

\subsection{Delimitación y validación de los mensajes}

Cada mensaje es un objeto JSON codificado en UTF-8 y terminado por un byte
de salto de línea. Los saltos de línea dentro del texto se escapan durante
la serialización JSON. Una lectura de TCP no se interpreta como un mensaje
completo: el receptor acumula bytes hasta encontrar el delimitador y
conserva los sobrantes para la siguiente lectura. La decodificación UTF-8
ocurre después de reconstruir la línea, por lo que un carácter puede llegar
fragmentado en varias lecturas.

Las solicitudes contienen los campos \texttt{sender}, \texttt{timestamp},
\texttt{message} y \texttt{message\_id}. Se exige texto no vacío para el
emisor y el identificador, texto para el mensaje y un entero no negativo
para el reloj; los valores booleanos no se aceptan como enteros. El límite
de cada objeto JSON serializado es de 1 MiB, sin contar el delimitador.
Este límite incluye los campos y los escapes del JSON, además de la carga
útil. Los tamaños experimentales de 100 y 10240 bytes corresponden a la
carga útil definida por el ejecutor, no al tamaño total transmitido.

\subsection{Acuses y reloj de Lamport}

El receptor actualiza su reloj al aceptar la solicitud y vuelve a
incrementarlo antes de enviar el acuse. Este contiene los campos comunes,
\texttt{ok=true}, un identificador formado con el sufijo \texttt{-ack}
y el campo \texttt{ack\_for}, que confirma el identificador original.
El cliente comprueba el emisor, ambos identificadores, el indicador de
éxito y que el reloj del acuse sea posterior al de su envío. Después
aplica la actualización de recepción a su propio reloj.

Cada proceso mantiene una instancia de \texttt{RelojLamport}. Un bloqueo
adicional agrupa el avance del reloj con la escritura del evento, para
evitar que los hilos registren dos valores en orden inverso. Las solicitudes
y los acuses generan eventos de envío y recepción. El siguiente fragmento
resume la operación local usada en \texttt{EstadoNodo.evento}:

\begin{verbatim}
with self._lock:
    tipo = "envio" if remoto is None else "recepcion"
    valor = (self.reloj.envio() if remoto is None
             else self.reloj.recepcion(remoto))
    if self.registro is not None:
        self.registro.escribir(identificador, tipo, valor)
    return valor
\end{verbatim}

\subsection{Latencia y manejo de errores}

Se conserva una conexión por destino y puerto durante la vida del cliente.
La primera conexión se establece antes de iniciar el intervalo de latencia.
La medición usa \texttt{time.perf\_counter()} desde el envío de la solicitud,
incluida su serialización, hasta recibir y validar el acuse. El resultado
se expresa en milisegundos. En las corridas de latencia se omite el argumento
\texttt{--eventos}; el reloj permanece activo.

El tiempo máximo predeterminado es de cinco segundos para la operación
completa, incluida la conexión inicial. Un plazo absoluto se comprueba
en cada lectura para que una respuesta que llega lentamente no prolongue
indefinidamente la espera. Se puede configurar mediante \texttt{--timeout}
con valores mayores que cero y de hasta diez segundos, inferiores a los
quince segundos de espera del ejecutor.

Un destino no disponible, un cierre prematuro, un JSON inválido o un acuse
incorrecto producen una respuesta con \texttt{ok=false} y una descripción
del error. Se descarta la conexión afectada; no se reintenta automáticamente
la solicitud, porque el receptor podría haberla procesado. Una orden
posterior puede abrir otra conexión. Los diagnósticos se escriben en la
salida de error para preservar el protocolo JSON de la salida estándar.

\subsection{Ejecución y comprobación de desarrollo}

Desde la raíz del repositorio se inicia cada nodo en una terminal distinta.
Por ejemplo, para A:

\begin{verbatim}
python -m src.sockets.nodo --nombre nodo-A --puerto 5000 ^
  --corrida prueba-khriz
\end{verbatim}

El carácter de continuación mostrado corresponde a \texttt{cmd.exe}.
Para B y C se sustituyen el nombre y el puerto por \texttt{nodo-B}/5001 y
\texttt{nodo-C}/5002. El nodo informa su disponibilidad mediante una línea
JSON. Una orden de ejemplo, escrita como una sola línea en la terminal de A,
es la siguiente:

\begin{verbatim}
{"destino":"nodo-B","puerto":5001,"message":"Hola","message_id":"K-001"}
\end{verbatim}

Para registrar eventos se añade \texttt{--eventos} seguido de una
ruta CSV distinta para cada nodo; el archivo debe ser nuevo. El proceso
cierra sus conexiones al terminar la entrada estándar.

La integración también se comprueba con el ejecutor existente:

\begin{verbatim}
python experiments/ejecutar.py --mecanismo sockets ^
  --modo humo --salida data/prueba-tcp-khriz
python -B -m unittest discover -s tests -v
\end{verbatim}

La carpeta de salida de la prueba de humo debe ser nueva. Las pruebas
automatizadas cubren reconstrucción de UTF-8 fragmentado, mensajes
agrupados, conexiones persistentes, mensajes largos, entradas inválidas,
tiempo de espera, acuses incorrectos y comunicación de los tres procesos.
Además verifican eventos reales de solicitudes y acuses y dos envíos
independientes con el mismo reloj. Sus CSV se crean en carpetas temporales
y no sustituyen los datos ni las conclusiones de la campaña experimental.

La captura de autoría del intercambio se guarda en
\texttt{evidencias/khriz/}. Debe mostrar la ejecución real en el equipo
del integrante, junto con su identificación y la fecha y hora del sistema.
\section{Implementación de la comunicación mediante gRPC}
\label{sec:grpc}

Cada nodo mantiene un servidor gRPC y un cliente dentro del mismo proceso.
El servidor atiende RPC mediante un grupo de cuatro hilos y escucha en
\texttt{127.0.0.1}. El cliente utiliza un canal por destino y puerto.
\texttt{nodo.py} integra ambos componentes con el reloj de Lamport y el
registro CSV de \texttt{src/comun/}, sin depender del transporte TCP.

\subsection{Contrato y generación}

El servicio \texttt{Mensajeria} expone la operación unaria
\texttt{EnviarMensaje}: cada solicitud produce un acuse. El contrato
de \texttt{proto/mensajes.proto} contiene los mismos datos lógicos que TCP:

\begin{verbatim}
syntax = "proto3";
package pe_u1;
service Mensajeria {
  rpc EnviarMensaje(Mensaje) returns (Acuse);
}
message Mensaje {
  string sender = 1;
  int64 timestamp = 2;
  string message = 3;
  string message_id = 4;
}
message Acuse {
  string sender = 1;
  int64 timestamp = 2;
  string message = 3;
  string message_id = 4;
  string ack_for = 5;
  bool ok = 6;
}
\end{verbatim}

Los campos de texto usan \texttt{string} y el reloj usa un entero de
64 bits con signo. El servicio rechaza relojes negativos e identificadores
o emisores vacíos mediante el estado \texttt{INVALID\_ARGUMENT}. Los números
de campo identifican los valores en la representación binaria. El acuse
confirma el identificador original mediante \texttt{ack\_for}; su propio
identificador añade el sufijo \texttt{-ack}.

La generación sigue el procedimiento documentado por
gRPC\footnote{\texttt{https://grpc.io/docs/languages/python/basics/}}.
Desde la raíz, en \texttt{cmd.exe}, se ejecuta:

\begin{verbatim}
python -m grpc_tools.protoc -Isrc/grpc/generado=proto ^
  --python_out=. --grpc_python_out=. proto/mensajes.proto
\end{verbatim}

El prefijo virtual determina la ruta \texttt{src/grpc/generado/} y las
importaciones del paquete. Allí se generan \texttt{mensajes\_pb2.py}
y \texttt{mensajes\_pb2\_grpc.py}, sin modificaciones manuales.

Se comprobaron las siguientes versiones: Python 3.12, grpcio 1.85.3,
grpcio-tools 1.85.3 y protobuf 7.36.2.

\subsection{Reloj, acuses y latencia}

Cada proceso comparte un único reloj entre cliente y servidor.
Antes de enviar la solicitud se incrementa el contador. El servidor
aplica \(\max(local, recibido)+1\) al recibirla y vuelve a incrementarlo
al enviar el acuse. El cliente valida el emisor, los identificadores,
el indicador de éxito y que el reloj del acuse sea posterior al de su
solicitud; después registra la recepción. Un bloqueo agrupa cada avance
del reloj con su escritura CSV para conservar el orden local entre hilos.

La conexión inicial se establece antes del intervalo de latencia y el
canal se reutiliza en las solicitudes posteriores. La medición incluye
la serialización de la solicitud, el RPC, la deserialización del acuse
y su validación, y termina antes del registro de la recepción:

\begin{verbatim}
inicio = perf_counter()
acuse = stub.EnviarMensaje(solicitud, timeout=restante,
                          wait_for_ready=False)
\end{verbatim}

Después de validar el acuse se calcula
\texttt{(perf\_counter() - inicio) * 1000}. El tiempo máximo de la
operación completa, incluida la conexión inicial, es de cinco segundos
por defecto. El RPC recibe el tiempo restante de ese presupuesto.
\texttt{--timeout} permite valores mayores que cero y de hasta diez
segundos, inferiores a la espera de quince segundos del ejecutor.
Tras un fallo se cierra el canal afectado. La aplicación no reintenta
automáticamente la solicitud, porque el servidor podría haberla procesado
antes de agotarse el plazo. Una orden posterior puede abrir otro canal.

\subsection{Interfaz de ejecución y errores}

Los argumentos, las órdenes JSON por entrada estándar y las respuestas
son compatibles con el ejecutor existente. Una respuesta inicial con
\texttt{ok=true} anuncia que el servidor está activo. Las órdenes contienen
\texttt{destino}, \texttt{puerto}, \texttt{message} y \texttt{message\_id}.
Una respuesta exitosa devuelve el identificador, \texttt{ok=true} y
\texttt{latencia\_ms}. Los fallos devuelven \texttt{ok=false} y una
descripción; los errores gRPC incluyen el estado y el detalle del RPC.
La salida estándar contiene exclusivamente respuestas JSON. Los
diagnósticos se reservan para la salida de error. Al terminar la entrada
estándar, el proceso cierra el cliente y el servidor.

Para iniciar A desde la raíz se utiliza:

\begin{verbatim}
python -m src.grpc.nodo --nombre nodo-A --puerto 5000 ^
  --corrida prueba-byron --eventos data/grpc/prueba-A.csv
\end{verbatim}

B y C se inician en terminales distintas con sus nombres y puertos
5001 y 5002, y con rutas CSV propias. Los archivos de eventos deben ser
nuevos. Cuando se omite \texttt{--eventos}, el reloj sigue activo pero
no se escriben eventos, conforme al protocolo de medición.

\subsection{Comprobación manual en el equipo de Byron}

El 21 de noviembre de 2026, Byron inició los tres nodos en terminales
de su equipo y ejecutó tres solicitudes: \texttt{B-001} de A hacia B,
\texttt{B-002} de B hacia C y \texttt{B-003} de C hacia A. Las tres
recibieron un acuse válido. Después cerró los procesos mediante EOF.
Las capturas originales de los intercambios y del cierre se conservan
en \texttt{evidencias/byron/}, con usuario, equipo, fecha y zona horaria
visibles.

Los tres CSV de \texttt{data/grpc/practica-byron-01/} contienen doce
eventos: envío y recepción de tres solicitudes y sus tres acuses.
El primer intercambio registra los valores 1, 2, 3 y 4, correspondientes
al envío de A, recepción en B, envío del acuse de B y recepción en A.
En los seis pares registrados, el reloj de recepción es mayor que el
de envío. Los relojes locales también crecen estrictamente.
Los eventos abarcan desde las 18:49:51 hasta las 18:52:27, UTC-05:00.
Después del cierre se calcularon y verificaron las sumas de los CSV
en el archivo \texttt{SHA256SUMS} de esa misma carpeta.

Esta comprobación manual documenta el intercambio entre los tres
procesos. Por su número de envíos y sus condiciones, no sustituye
las repeticiones, tamaños y concurrencia de la campaña experimental,
ni permite concluir qué transporte tiene menor latencia.

\subsection{Asistencia en el desarrollo}

La declaración individual del bloque gRPC se recoge en la
sección~\ref{sec:declaracion-ia}. La bitácora y los archivos de la práctica
conservan el detalle de las intervenciones y de las comprobaciones.

\section{Protocolo experimental}
Se programan veinte solicitudes para comprobar Lamport. Una corrida adicional
parte de relojes reiniciados y programa los primeros envíos de A y B hacia C.
Para medir latencia se usan cargas útiles de 100 y 10240 bytes, modos
secuencial y concurrente y cien solicitudes por corrida. Cada combinación
tiene un calentamiento y cinco repeticiones válidas. La latencia mide el
intervalo entre el envío y el acuse. El rendimiento se calcula con las
solicitudes completadas y la duración de pared de cada corrida.
% Incorporar las condiciones reales a partir de entorno.md.

\section{Resultados y discusión}
Pendientes de la campaña oficial y del análisis de John.

\section{Aportes}
% Completar con aportes efectivamente realizados y enlaces de los cuatro integrantes.

\section{Declaración de uso de inteligencia artificial}
\label{sec:declaracion-ia}

\noindent
\begin{tabularx}{\linewidth}{|>{\raggedright\arraybackslash}p{0.17\linewidth}|>{\raggedright\arraybackslash}p{0.21\linewidth}|>{\raggedright\arraybackslash}X|}
\hline
\textbf{Integrante} & \textbf{Herramienta} & \textbf{Uso declarado} \\
\hline
Khriz Coronel & Pendiente & Pendiente \\
\hline
Diego Giraldo & Codex (OpenAI) &
Apoyo en la preparación de la base del ejecutor, la redacción y el formato del informe, y la búsqueda y verificación de referencias. \\
\hline
Byron Daniel Loor Mendoza & Codex (OpenAI), agente basado en GPT-6 &
Utilicé OpenAI Codex como herramienta de asistencia para organizar el
repositorio y desarrollar el bloque gRPC. Su intervención comprendió la
elaboración del contrato de mensajes, la implementación, la documentación
técnica y la revisión de su presentación. También realizó comprobaciones
automatizadas de desarrollo y verificaciones de los registros y las sumas
SHA-256. El detalle de la asistencia, los archivos afectados y las
comprobaciones realizadas está documentado en la bitácora y el historial
del repositorio. \\
\hline
John Silva & Pendiente & Pendiente \\
\hline
\end{tabularx}


\printbibliography
\end{document}
